\documentclass[11pt,a4paper]{article}
\usepackage[utf8]{inputenc}
\usepackage[spanish]{babel}
\usepackage{amsmath, amssymb}
\usepackage{graphicx}
\usepackage{booktabs}
\usepackage{geometry}
\usepackage{float}
\usepackage{hyperref}
\usepackage{listings}
\usepackage{xcolor}

\geometry{a4paper, margin=2.54cm}

\lstset{
    language=Python,
    basicstyle=\ttfamily\small,
    keywordstyle=\color{blue},
    stringstyle=\color{red},
    commentstyle=\color{green!50!black},
    showstringspaces=false,
    frame=single,
    breaklines=true
}

\begin{document}

\begin{titlepage}
    \vspace*{-2cm}
    \begin{center}
        \includegraphics[width=3cm]{logoUANL.png} \hspace{2cm}
        \includegraphics[width=3cm]{logoFCFM.png}
    \end{center}
    \vspace{2cm}
    \begin{center}
        \Huge \textbf{Producto Integrador de Aprendizaje (PIA)}\\[0.5cm]
        \LARGE Investigación de Operaciones (IO)\\[0.5cm]
        \Large Facultad de Ciencias Físico Matemáticas, UANL\\
        Semestre Enero-Junio 2026\\[3cm]
        
        \Large
        \textbf{Nombre:} Diego Andre Islas Cadillo\\
        \textbf{Matrícula:} 2132810\\[1cm]
        
        \textbf{Tema:}\\
        8. Proyecto CPM para el mantenimiento preventivo mayor \\
        (paro de planta) de una refinería.\\[2cm]
        
        \rule{0.8\textwidth}{0.5pt}\\[0.5cm]
        \normalsize \textbf{Declaración de Originalidad:}\\
        \textit{Declaro que este código es de mi autoría y ha sido desarrollado bajo los criterios de integridad académica de la FCFM.}\\[0.5cm]
        \rule{0.8\textwidth}{0.5pt}\\[2cm]
        
        \large \today
    \end{center}
\end{titlepage}

\section{Introducción y Contexto del Problema}

El presente proyecto aborda un desafío crítico en la industria de los hidrocarburos: el mantenimiento preventivo mayor (también conocido como \textit{Turnaround} o Paro de Planta). Este escenario se centra en la \textbf{Unidad de Destilación Atmosférica} de la empresa ficticia \textbf{Refinería ``PetroNorte''}.

Un paro de planta implica detener por completo la producción para realizar inspecciones profundas, reparaciones, limpieza y reemplazo de equipos críticos (como torres de destilación, intercambiadores de calor y bombas). Debido a que cada día que la planta está detenida representa pérdidas millonarias en producción no realizada, el tiempo es el recurso más crítico.

\paragraph{Problemática a resolver:}
La gerencia de operaciones de PetroNorte necesita determinar el tiempo mínimo posible para completar el \textit{Turnaround} de este año, identificar las actividades críticas que no pueden retrasarse bajo ninguna circunstancia, y establecer los márgenes de holgura para las actividades no críticas. Esto permitirá una asignación óptima de los contratistas y recursos disponibles.

\section{Definición de Datos y Variables}

Para modelar este escenario, se han definido 20 actividades macro que representan las tareas estandarizadas en un paro de planta de esta magnitud. Los tiempos (en horas) han sido estimados considerando cuadrillas trabajando en turnos continuos de 24 horas.

\paragraph{Justificación de realismo:}
\begin{itemize}
    \item Las dependencias lógicas respetan el protocolo industrial de seguridad (e.g., el bloqueo eléctrico ``LOTO'' antes de abrir equipos).
    \item Las duraciones de pruebas (como las Pruebas No Destructivas y las Pruebas Hidrostáticas) reflejan estándares normativos de inspección API (\textit{American Petroleum Institute}).
\end{itemize}

\begin{table}[H]
\centering
\begin{tabular}{cllc}
\toprule
\textbf{ID} & \textbf{Actividad} & \textbf{Duración (h)} & \textbf{Predecesoras} \\
\midrule
A & Paro de unidad y despresurización & 36 & - \\
B & Drenado, purga y lavado con vapor & 72 & A \\
C & Aislamiento eléctrico y bloqueo de válvulas (LOTO) & 24 & B \\
D & Apertura de registros de la torre de destilación & 48 & C \\
E & Desmontaje e inspección de bombas de carga & 60 & C \\
F & Extracción de haces de tubos en intercambiadores & 72 & C \\
G & Limpieza e inspección de platos de la torre & 108 & D \\
H & Reemplazo de impulsores y sellos en bombas & 48 & E \\
I & Limpieza con hidrojet de intercambiadores & 72 & F \\
J & Pruebas No Destructivas (NDT) en tuberías & 36 & D, F \\
K & Reparación de soldaduras en tuberías & 90 & J \\
L & Reinstalación de haces de tubos & 54 & I, K \\
M & Cierre de registros de la torre & 36 & G, J \\
N & Montaje y alineación de bombas & 48 & H \\
O & Pruebas hidrostáticas del sistema & 72 & L, M, N \\
P & Retiro de bloqueos (LOTO) y energización & 24 & O \\
Q & Calibración y prueba de lazos de control & 48 & O \\
R & Restablecimiento de aislamiento térmico & 60 & O \\
S & Recirculación en frío & 36 & P, Q \\
T & Arranque y estabilización de planta & 72 & R, S \\
\bottomrule
\end{tabular}
\caption{Tabla de Actividades del Proyecto CPM}
\end{table}

\section{Modelado Matemático (CPM)}

El Método de la Ruta Crítica (CPM) utiliza las siguientes formulaciones matemáticas en un grafo dirigido acíclico $G = (V, E)$:

\begin{enumerate}
    \item \textbf{Early Start ($ES_i$) y Early Finish ($EF_i$):}
    \begin{align*}
        ES_i &= \max_{j \in Pred(i)} (EF_j) \\
        EF_i &= ES_i + T_i
    \end{align*}
    Donde $T_i$ es la duración de la actividad $i$. Para el nodo de inicio, $ES_{\text{inicio}} = 0$.

    \item \textbf{Late Start ($LS_i$) y Late Finish ($LF_i$):}
    \begin{align*}
        LF_i &= \min_{j \in Succ(i)} (LS_j) \\
        LS_i &= LF_i - T_i
    \end{align*}
    Para el nodo final, $LF_{\text{fin}} = EF_{\text{fin}}$ (Duración total del proyecto).

    \item \textbf{Holgura ($S_i$):}
    \begin{equation*}
        S_i = LS_i - ES_i
    \end{equation*}
\end{enumerate}

\textbf{Condición de Ruta Crítica:} Una actividad pertenece a la ruta crítica si y solo si su holgura es cero ($S_i = 0$).

\section{Desarrollo Computacional (Python)}

Para la solución computacional del modelo CPM, se implementó un script en Python utilizando las siguientes librerías profesionales:
\begin{itemize}
    \item \textbf{NetworkX:} Fundamental para la creación del grafo dirigido (\texttt{nx.DiGraph}) y la resolución de la topología de red para evaluar los tiempos tempranos (Forward Pass) y tiempos tardíos (Backward Pass).
    \item \textbf{Pandas:} Utilizada para la estructuración de la tabla de datos inicial y la tabulación final de los resultados, permitiendo calcular fácilmente la ruta crítica y las holguras de cada actividad.
    \item \textbf{Matplotlib:} Empleada para el renderizado del Diagrama de Gantt automatizado, logrando una visualización gerencial clara de las actividades críticas.
\end{itemize}

A continuación, se presentan los fragmentos clave del código fuente que resuelven matemáticamente el problema:

\subsection*{Fragmento 1: Construcción de la Red Topológica}
\begin{lstlisting}
import networkx as nx
import pandas as pd

# Creacion del Grafo Dirigido
G = nx.DiGraph()

# Anadir nodos y duraciones
for _, row in df_actividades.iterrows():
    G.add_node(row["Actividad"], duration=row["Duración"])

# Definicion de dependencias logicas
for _, row in df_actividades.iterrows():
    for pred in row["Predecesoras"]:
        G.add_edge(pred, row["Actividad"])
\end{lstlisting}

\subsection*{Fragmento 2: Cálculo de la Ruta Crítica (Forward y Backward Pass)}
\begin{lstlisting}
# Calculo de ES y EF (Forward Pass)
for node in nx.topological_sort(G):
    es = max([G.nodes[pred]["EF"] for pred in G.predecessors(node)], default=0)
    G.nodes[node]["ES"] = es
    G.nodes[node]["EF"] = es + G.nodes[node]["duration"]

# Calculo de LS y LF (Backward Pass)
project_duration = G.nodes["Fin"]["EF"]
for node in reversed(list(nx.topological_sort(G))):
    lf = min([G.nodes[succ]["LS"] for succ in G.successors(node)], default=project_duration)
    G.nodes[node]["LF"] = lf
    G.nodes[node]["LS"] = lf - G.nodes[node]["duration"]
    
# Calculo de Holgura (Slack)
for node in G.nodes():
    G.nodes[node]["Slack"] = G.nodes[node]["LS"] - G.nodes[node]["ES"]
\end{lstlisting}

\section{Conclusiones Ejecutivas}

Tras el análisis computacional del modelo CPM, se derivan las siguientes conclusiones operativas para la toma de decisiones gerenciales en la Refinería PetroNorte:

\begin{figure}[H]
    \centering
    \includegraphics[width=1\textwidth]{gantt_chart.png}
    \caption{Diagrama de Gantt del Proyecto y Ruta Crítica.}
\end{figure}

\begin{enumerate}
    \item \textbf{Duración Mínima del Proyecto:} El paro de planta tomará un mínimo de \textbf{612 horas continuas (25.5 días)}. Cualquier compromiso de arrancar la unidad antes de esta fecha es matemáticamente imposible con los recursos y tiempos actuales.
    \item \textbf{Identificación de la Ruta Crítica:} La ruta crítica está conformada por la secuencia: \\
    \texttt{A $\rightarrow$ B $\rightarrow$ C $\rightarrow$ F $\rightarrow$ J $\rightarrow$ K $\rightarrow$ L $\rightarrow$ O $\rightarrow$ Q $\rightarrow$ S $\rightarrow$ T}.
    \begin{itemize}
        \item \textbf{Interpretación Gerencial:} Las actividades de "Extracción de tubos", "Pruebas no destructivas", "Reparación de soldaduras" y "Pruebas hidrostáticas" no tienen margen de error. Si las reparaciones de tuberías (Actividad K) se retrasan un solo día, el arranque de toda la planta se retrasará exactamente un día.
    \end{itemize}
    \item \textbf{Gestión de Holguras:} Existen tareas como el mantenimiento de bombas (E, H, N) y limpieza de platos de torres (G) que cuentan con holguras considerables (hasta 96 horas en el caso de las bombas). 
    \begin{itemize}
        \item \textbf{Acción Sugerida:} La gerencia puede desviar personal temporalmente de las cuadrillas de bombas para apoyar en las actividades de intercambiadores y soldaduras (ruta crítica) sin afectar la fecha final de entrega del proyecto.
    \end{itemize}
    \item \textbf{Optimización de Presupuesto:} Dado que los contratistas de instrumentación tienen una holgura moderada en algunas tareas iniciales, sus horas extra solo deben ser autorizadas si están trabajando en las actividades O, Q, S y T.
\end{enumerate}

El modelo en Python automatiza exitosamente el control logístico, permitiendo recalcular la ruta crítica de inmediato si ocurre un imprevisto durante el mantenimiento real.

\end{document}
